\subsection{分次}

设 \(\Delta\) 是一个以加法记的交换幺半群。设 \(\phi_0\) 表示 X 到 \(\Delta\) 中的一个映射，并用 \(\phi\) 表示自由 magma M(X) 到 \(\Delta\) 中的扩张 \(\phi_0\) 的同态。对一切 \(\delta \in \Delta\) ，令 \(\operatorname{Lib}^{\delta}(X)\) 为 \(\operatorname{Lib}(X)\) 中以 M(X) 的子集 \(\phi^{-1}(\delta)\) 为基的子模。族 \((\operatorname{Lib}^{\delta}(X))_{\delta \in \Delta}\) 是代数 \(\operatorname{Lib}(X)\) 的一个分次，即

\begin{enumerate}
\def\labelenumi{(\arabic{enumi})}
\setcounter{enumi}{7}
\tightlist
\item
\end{enumerate}

\[
\operatorname{Lib} (\mathrm{X}) = \bigoplus_ {\delta \in \Delta} \operatorname{Lib} ^ {\delta} (\mathrm{X})\tag{8}
\]

\[
\operatorname{Lib} ^ {\delta} (\mathrm{X}). \operatorname{Lib} ^ {\delta^ {\prime}} (\mathrm{X}) \subset \operatorname{Lib} ^ {\delta + \delta^ {\prime}} (\mathrm{X}) \quad \text { for } \delta , \delta^ {\prime} \text { in } \Delta\tag{9}
\]

(《代数》， 第 III 章， § 3， no. 1， 例 3)。

引理 1. 定义 1 中的理想 \(\mathfrak{a}\) 是分次的。

设 \(a, b\) 属于 \(\operatorname{Lib}(\mathbf{X})\) ，令 \(\mathbf{B}(a, b) = a.b + b.a\) 。公式

\begin{enumerate}
\def\labelenumi{(\arabic{enumi})}
\setcounter{enumi}{9}
\tightlist
\item
\end{enumerate}

\[
\mathrm{B} (a, b) = \mathrm{Q} (a + b) - \mathrm{Q} (a) - \mathrm{Q} (b)\tag{10}
\]

\[
Q \left(\lambda_ {1}. w _ {1} + \dots + \lambda_ {n}. w _ {n}\right) = \sum_ {i} \lambda_ {i} ^ {2} Q \left(w _ {i}\right) + \sum_ {i <   j} \lambda_ {i} \lambda_ {j} B \left(w _ {i}, w _ {j}\right)\tag{11}
\]

对 M(X) 中的 \(w_1, \ldots, w_n\) 与 K 中的 \(\lambda_1, \ldots, \lambda_n\) 表明，族 \((\mathbf{Q}(a))_{a \in \mathrm{Lib}(\mathbf{X})}\) 与 \((\mathbf{Q}(w), \mathbf{B}(w, w'))_{w, w' \in \mathrm{M}(\mathbf{X})}\) 生成 Lib(X) 的同一个子模。由于 J 是三线性的，理想 a 由齐次元素 \(\mathbf{Q}(w), \mathbf{B}(w, w')\) 与 \(\mathbf{J}(w, w', w'')\) 生成，其中 \(w, w', w''\) 属于 M(X)，因而它是分次的 (《代数》， 第 III 章， § 3， no. 3， 命题 1)。

给李代数 \(\mathrm{L}(\mathbf{X}) = \mathrm{Lib}(\mathbf{X}) / a\) 赋予商分次。\(\mathrm{L}(\mathbf{X})\) 的 \(\delta\) 次齐次分量记作 \(\mathrm{L}^{\delta}(\mathbf{X})\) ；它是 \(\mathrm{L}(\mathbf{X})\) 中由元素 \(w \in \mathrm{M}(\mathbf{X})\) （它们满足 \(\phi(w) = \delta\) ）的像生成的子模。

我们将特别用到以下两种情形：

\begin{enumerate}
\def\labelenumi{(\alph{enumi})}
\tightlist
\item
  全分次：我们取 \(\Delta = \mathbf{N}\) ，\(\phi_0(x) = 1\) （对一切 \(x\in \mathbf{X}\) ），从而得到 \(\phi (w) = l(w)\) ，其中 \(w\) 属于 \(\mathrm{M}(\mathrm{X})\) 。K-模 \(\mathrm{L}^n (\mathrm{X})\) 由长度为 \(n\) 、属于 \(\mathrm{M}(\mathrm{X})\) 的元素的像生成，我们称这些元素为 \(n\) 次交错元。后面将会看到，模 \(\mathrm{L}^n (\mathrm{X})\) 是自由的，并且具有由 \(n\) 次交错元组成的一个基 (no. 11， 定理 1)。于是 \(\mathrm{L}(\mathrm{X}) = \bigoplus_{n\geq 1}\mathrm{L}^n (\mathrm{X})\) ，并且 \(\mathrm{L}^1 (\mathrm{X})\) 以 X 为基 (no. 2， 命题 1 的推论 1)。按照 \(\mathrm{M}(\mathrm{X})\) 的构造，
\end{enumerate}

\[
\mathrm{L} ^ {n} (\mathrm{X}) = \sum_ {p = 1} ^ {n - 1} [ \mathrm{L} ^ {p} (\mathrm{X}), \mathrm{L} ^ {n - p} (\mathrm{X}) ]\tag{12}
\]

特别地

\[
[ \mathrm{L} ^ {m} (\mathrm{X}), \mathrm{L} ^ {n} (\mathrm{X}) ] \subset \mathrm{L} ^ {m + n} (\mathrm{X}).\tag{13}
\]

\begin{enumerate}
\def\labelenumi{(\alph{enumi})}
\setcounter{enumi}{1}
\tightlist
\item
  多重分次：我们取 \(\Delta\) 为在 X 上构造的自由交换幺半群 \(\mathbf{N}^{(\mathbf{X})}\) 。映射 \(\phi_0\) 把 X 映到 \(\Delta\) 中，由 \((\phi_0(x))(x') = \delta_{xx'}\) 定义，其中 \(\delta_{xx'}\) 是 Kronecker 符号。对 \(w \in M(X)\) 与 \(x \in X\) ，整数 \((\phi(w))(x)\) 是“字母 \(x\) 在 \(w\) 中出现的次数”。设 \(\alpha\) 属于 \(\mathbf{N}^{(\mathbf{X})}\) ，我们记 \(|\alpha| = \sum_{x \in X} \alpha(x)\) ，从而 \(|\phi(w)| = l(w)\) （对 M(X) 中的一切 \(w\) ）。由此可得
\end{enumerate}

\[
\mathbf {L} ^ {n} (\mathbf {X}) = \bigoplus_ {| \alpha | = n} \mathbf {L} ^ {\alpha} (\mathbf {X});\tag{14}
\]

显然

\[
[ \mathrm{L} ^ {\alpha} (\mathrm{X}), \mathrm{L} ^ {\beta} (\mathrm{X}) ] \subset \mathrm{L} ^ {\alpha + \beta} (\mathrm{X}) \quad \text { for } \alpha , \beta \text { in } \mathbf {N} ^ {(\mathrm{X})}.\tag{15}
\]

命题 4. 设 \(S\) 是 \(X\) 的一个子集。若将 \(\mathbf{N}^{(\mathbb{S})}\) 与它在 \(\mathbf{N}^{(\mathbf{X})}\) 中的典范像等同 (《代数》， 第 I 章， § 7， no. 7)，则 \(L(S) = \sum_{\alpha \in \mathbf{N}^{(\mathbb{S})}} L^{\alpha}(X)\) 。此外，对一切 \(\alpha \in \mathbf{N}^{(\mathbb{S})}\) ， \(\alpha\) 次齐次分量（在 \(L(S)\) 上的多重分次之下）等于 \(L^{\alpha}(X)\) 。

 设 \(\alpha \in \mathbf{N}^{(S)}\) 。模 \(\mathrm{L}^{\alpha}(\mathrm{S})\) 由 \(\mathrm{L}(\mathrm{X})\) 中的像生成，即元素 \(w\) （属于 \(\mathrm{M}(\mathrm{S})\) 且满足 \(\phi(w) = \alpha\) ）的像，也就是说 (《代数》， § 7， no. 9， 公式 (23) 与 (24))，它是 \(w\) 组成的集合，其中元素属于 \(\mathrm{M}(\mathrm{X})\) 且满足 \(\phi(w) = \alpha\) 。由此得 \(\mathrm{L}^{\alpha}(\mathrm{S}) = \mathrm{L}^{\alpha}(\mathrm{X})\) 。命题由这一点以及关系式 \(\mathrm{L}(\mathrm{S}) = \sum_{\alpha \in \mathbf{N}^{(S)}}\mathrm{L}^{\alpha}(\mathrm{S})\) 得证。

推论. 对每个子集族 \((\mathbf{S}_i)_{i\in I}\) （子集取自 \(\mathbf{X}\) ），

\[
\mathrm{L} \left(\bigcap_ {i \in I} S _ {i}\right) = \bigcap_ {i \in I} \mathrm{L} (S _ {i}).\tag{16}
\]

这由命题 4 以及显然的公式

\[
\mathbf {N} ^ {(S)} = \bigcap_ {i \in I} \mathbf {N} ^ {(S _ {i})}\tag{17}
\]

得出，其中我们已记 \(S = \bigcap_{i \in I} S_i\) 。

